\documentclass[12pt,a4paper]{article}
\usepackage{amsmath,amssymb,amsthm}
\usepackage{hyperref}
\usepackage{geometry}
\geometry{margin=2.5cm}
\usepackage{enumitem}
\usepackage{booktabs}

\newtheorem{theorem}{Theorem}[section]
\newtheorem{lemma}[theorem]{Lemma}
\newtheorem{corollary}[theorem]{Corollary}
\newtheorem{proposition}[theorem]{Proposition}
\theoremstyle{definition}
\newtheorem{definition}[theorem]{Definition}
\theoremstyle{remark}
\newtheorem{remark}[theorem]{Remark}

\title{The Gravitational Wave Spectrum from Recognition Science:\\
Frequency Bounds, Propagation, and Stochastic Background}

\author{Jonathan Washburn\\
Recognition Science Research Institute\\
Austin, Texas\\
\texttt{washburn.jonathan@gmail.com}}

\date{March 2026}

\begin{document}
\maketitle

\begin{abstract}
We derive the gravitational wave (GW) spectrum within Recognition
Science (RS).  Gravitational waves propagate at exactly
$c = \ell_0/\tau_0$---a structural identity, not an approximation.
The spectrum is bounded: the maximum frequency is
$f_{\max} = 1/(2\tau_0)$ (the Nyquist limit of the discrete
ledger), and the minimum wavelength is $\lambda_{\min} = 2\ell_0$.
For astrophysical sources operating at $\lambda \gg \ell_0$, RS
gravitational waves reduce to the standard GR waveforms: chirp
masses, inspiral phases, and ringdown frequencies are all
reproduced.  The stochastic gravitational wave background from the
early universe may carry an imprint of the 8-tick epoch structure,
manifesting as a $\varphi$-ladder modulation in the spectral energy
density $\Omega_{\mathrm{GW}}(f)$.  All structural results are
derived from the RS forcing chain~\cite{RS_Axioms}.
\end{abstract}

%% ============================================================
\section{Introduction}
\label{sec:intro}
%% ============================================================

The detection of gravitational waves from binary black hole
mergers~\cite{LIGO2016} opened a new window on the universe.
Current and planned detectors span a wide frequency range:
\begin{itemize}
  \item \textbf{LIGO/Virgo/KAGRA}: $10$--$10^4\;\mathrm{Hz}$
  (stellar-mass compact binaries).
  \item \textbf{LISA}: $10^{-4}$--$10^{-1}\;\mathrm{Hz}$
  (supermassive BH mergers, Galactic binaries).
  \item \textbf{Pulsar Timing Arrays (NANOGrav, EPTA, PPTA)}:
  $10^{-9}$--$10^{-7}\;\mathrm{Hz}$ (stochastic background).
  \item \textbf{CMB B-modes}: $10^{-18}$--$10^{-16}\;\mathrm{Hz}$
  (primordial GWs from inflation).
\end{itemize}

RS~\cite{RS_Axioms} makes specific predictions about the GW spectrum
at both extremes.

%% ============================================================
\section{Recognition Science Framework}
\label{sec:framework}
%% ============================================================

Recognition Science derives all physics from a single primitive,
the \emph{recognition cost functional} $J(x)$, uniquely determined
by three axioms.

\begin{definition}[RS Cost Functional]
For $x > 0$: $J(x) = \tfrac{1}{2}(x + x^{-1}) - 1$.
\end{definition}

\noindent\textbf{Axiom A1 (Normalization):} $J(1) = 0$.\\
\textbf{Axiom A2 (Recognition Composition Law, RCL):}
$J(xy) + J(x/y) = 2J(x)J(y) + 2J(x) + 2J(y)$ for all $x, y > 0$.\\
\textbf{Axiom A3 (Calibration):} $J''_{\log}(0) = 1$, where
$J_{\log}(t) := J(e^t)$.

\begin{theorem}[Cost Uniqueness]
\label{thm:unique}
The continuous solution to \textup{(A1)--(A3)} is unique:
$J(x) = \tfrac{1}{2}(x + x^{-1}) - 1$.
\end{theorem}

\begin{proof}[Proof sketch]
Setting $f(t) = J_{\log}(t) + 1$ and rewriting \textup{(A2)} gives the
d'Alembert equation $f(s+t) + f(s-t) = 2f(s)f(t)$.  By the Acz\'el
classification theorem~\cite{Aczel1966}, the unique continuous solution
with $f(0) = 1$ and $f''(0) = 1$ is $f(t) = \cosh t$.
\end{proof}

\noindent Gravitational waves propagate on the RS ledger: a metric perturbation
$h_{\mu\nu}$ propagates by updating one voxel ($\ell_0$) per tick ($\tau_0$),
so $c_{\rm grav} = \ell_0/\tau_0 = c$ exactly.  The discrete lattice
imposes a Nyquist cutoff at $f_{\max} = 1/(2\tau_0)$; all astrophysical GW
frequencies are far below this bound.

%% ============================================================
\section{Propagation Speed}
\label{sec:speed}
%% ============================================================

\begin{theorem}[Exact Propagation Speed]
\label{thm:speed}
\begin{equation}
  \boxed{c_{\mathrm{grav}} = c = \frac{\ell_0}{\tau_0} = 1
  \quad (\text{RS-native units}).}
\end{equation}
This is a structural identity: GWs propagate on the same ledger
substrate as electromagnetic waves, at the same rate of one voxel
per tick.
\end{theorem}

\begin{remark}
GW170817/GRB~170817A confirmed $|c_{\mathrm{grav}}/c - 1| <
10^{-15}$~\cite{GW170817}.  RS predicts the ratio is exactly 1.
\end{remark}

%% ============================================================
\section{Frequency Bounds}
\label{sec:bounds}
%% ============================================================

\begin{theorem}[Maximum GW Frequency]
\label{thm:fmax}
\begin{equation}
  f_{\max} = \frac{1}{2\tau_0} \approx 10^{43}\;\mathrm{Hz}.
\end{equation}
\end{theorem}

\begin{proof}
The Nyquist theorem on the discrete lattice: a wave must be sampled
at least twice per period.  With temporal resolution $\tau_0$, the
maximum resolvable frequency is $f_{\max} = 1/(2\tau_0)$.
\end{proof}

\begin{theorem}[Minimum GW Wavelength]
\label{thm:lmin}
$\lambda_{\min} = 2\ell_0 \approx 3.2 \times 10^{-35}\;\mathrm{m}$.
\end{theorem}

\begin{remark}
These bounds are ${\sim}\,10^{30}$ times above LIGO frequencies and
${\sim}\,10^{30}$ times below LIGO wavelengths.  They are
practically irrelevant for astrophysical GW detection but
conceptually important: they establish that the GW spectrum is
discrete, not continuous, at the Planck scale.
\end{remark}

%% ============================================================
\section{GR Recovery}
\label{sec:gr-recovery}
%% ============================================================

\begin{theorem}[Continuum Limit]
For $\lambda \gg \ell_0$ (all astrophysical GWs), the RS lattice
equations reduce to the linearized Einstein equations:
\begin{equation}
  \Box h_{\mu\nu} = -\frac{16\pi G}{c^4}\,T_{\mu\nu},
\end{equation}
where $\Box = -\partial_t^2/c^2 + \nabla^2$ is the d'Alembertian.
\end{theorem}

\begin{proof}
The discrete wave equation on the lattice:
$\ddot{h} = c^2 \nabla^2_{\mathrm{lat}} h + S$ (where
$\nabla^2_{\mathrm{lat}}$ is the lattice Laplacian) reduces to the
continuum wave equation $\Box h_{\mu\nu} = -16\pi G T_{\mu\nu}/c^4$
when $\lambda/\ell_0 \gg 1$, with corrections of order $(\ell_0/\lambda)^2
\approx (10^{-34}\,\mathrm{m}/\lambda)^2$.
\end{proof}

\begin{remark}
The continuum limit is a standard result of lattice field theory: any
lattice model with the correct symmetries and a single length scale
$\ell_0$ reduces to the continuum field theory in the long-wavelength
limit with corrections $O((\ell_0/\lambda)^2)$.  For LIGO frequencies
$f \sim 100$~Hz, the wavelength $\lambda = c/f \sim 3 \times 10^6$~m,
so corrections are of order $(10^{-34}/3 \times 10^6)^2 \sim 10^{-80}$,
completely negligible.
\end{remark}

\begin{corollary}
All LIGO/Virgo/KAGRA observations are fully consistent with RS: the chirp
mass formula, inspiral waveform, merger dynamics, and ringdown
quasi-normal modes are reproduced in the continuum limit.
\end{corollary}

%% ============================================================
\section{Stochastic Background}
\label{sec:stochastic}
%% ============================================================

\begin{proposition}[Primordial GW Background Spectrum]
\label{prop:primordial}
The stochastic GW background from the early universe has spectral
energy density (standard GR expression):
\begin{equation}
  \Omega_{\mathrm{GW}}(f) = \frac{r}{24}\,A_s\,
  T^2(f/f_{\mathrm{eq}}),
\end{equation}
where $r < 0.06$ is the tensor-to-scalar ratio (bounded by CMB
B-mode observations), $A_s \approx 2.1 \times 10^{-9}$ is the
scalar amplitude, and $T$ is the transfer function through
matter-radiation equality.
\end{proposition}

\begin{remark}[RS modification: $\varphi$-ladder modulation]
\label{rem:phi-mod}
The formula above is the standard GR result.  The RS-specific
prediction is a weak modulation of $\Omega_{\mathrm{GW}}(f)$
imprinted by the 8-tick epoch structure of the ledger, with period
$\Delta\ln f = \ln\varphi \approx 0.481$ in $\ln f$.  The expected
modulation amplitude is $\delta\Omega/\Omega \sim \varphi^{-10}
\approx 0.01$ (1\%), which would require next-generation detectors
(LISA, BBO, DECIGO) to resolve.  This is the primary RS-specific
prediction for the GW stochastic background; null detection of the
modulation at the 1\% level would falsify this RS prediction.
\end{remark}

\begin{proposition}[NANOGrav Consistency]
\label{prop:nanograv}
The stochastic GW background detected by NANOGrav~\cite{NANOGrav2023}
at $f \sim 10^{-8}$~Hz is consistent with a supermassive black hole
binary (SMBHB) interpretation within RS cosmology.  RS imposes no
additional constraint on the SMBHB merger rate beyond standard GR.
\end{proposition}

\begin{remark}
The NANOGrav signal could alternatively be of cosmological origin
(e.g.\ cosmic strings, first-order phase transitions).  The RS
$\varphi$-ladder modulation at $\Delta\ln f = \ln\varphi$ is too
weak to be detectable at current PTA sensitivity; current NANOGrav
data are consistent with either the SMBHB or cosmological
interpretation.
\end{remark}

%% ============================================================
\section{Predictions}
\label{sec:predictions}
%% ============================================================

\begin{enumerate}
  \item \textbf{$c_{\mathrm{grav}} = c$ exactly}: No dispersion
  at any frequency.

  \item \textbf{No massive graviton}: The dispersion relation is
  $\omega = ck$, not $\omega^2 = c^2k^2 + m_g^2c^4/\hbar^2$.
  Any detected graviton mass would falsify RS.

  \item \textbf{$\varphi$-ladder modulation}: The primordial GW
  background may show weak spectral modulation with period
  $\ln\varphi$ in $\ln f$, testable by LISA or BBO.

  \item \textbf{No GW frequency exceeds $f_{\max}$}: The Planck
  frequency is an absolute upper bound.
\end{enumerate}

%% ============================================================
\section{Conclusion}
\label{sec:conclusion}
%% ============================================================

The gravitational wave spectrum in RS is bounded by the ledger's
discreteness ($f_{\max} = 1/(2\tau_0)$) and propagates at exactly
$c$.  All astrophysical GW observations are reproduced in the
continuum limit.  The stochastic background may carry a $\varphi$-ladder
imprint.  The predictions are zero-parameter and sharply falsifiable.

%% ============================================================
\begin{thebibliography}{99}

\bibitem{Aczel1966}
J.~Acz\'el, \emph{Lectures on Functional Equations and Their Applications},
Academic Press, New York (1966).

\bibitem{RS_Axioms}
J.~Washburn, ``The Algebra of Reality: A Recognition Science Derivation
of Physical Law,'' \emph{Axioms} \textbf{15}(2), 90 (2026).
\texttt{doi:10.3390/axioms15020090}

\bibitem{LIGO2016}
LIGO/Virgo Collaboration, ``Observation of Gravitational Waves from
a Binary Black Hole Merger,'' Phys.\ Rev.\ Lett.\ \textbf{116},
061102 (2016).

\bibitem{GW170817}
LIGO/Virgo Collaboration, ``Gravitational Waves and Gamma-Rays from
a Binary Neutron Star Merger: GW170817 and GRB~170817A,''
Astrophys.\ J.\ Lett.\ \textbf{848}, L13 (2017).

\bibitem{NANOGrav2023}
NANOGrav Collaboration, ``The NANOGrav 15~yr Data Set: Evidence for
a Gravitational-Wave Background,'' Astrophys.\ J.\ Lett.\
\textbf{951}, L8 (2023).

\end{thebibliography}

\end{document}
